\documentclass[pdflatex,sn-mathphys-num]{sn-jnl}

\usepackage{graphicx}
\usepackage{multirow}
\usepackage{amsmath,amssymb,amsfonts}
\usepackage{amsthm}
\usepackage[title]{appendix}
\usepackage{xcolor}
\usepackage{textcomp}
\usepackage{manyfoot}
\usepackage{booktabs}
\usepackage{algorithm}
\usepackage{algorithmicx}
\usepackage{algpseudocode}
\usepackage{listings}

\graphicspath{{figs/}{./figs/}{./}}
\setkeys{Gin}{draft=false}

\raggedbottom

% The class loads article with twoside, so odd and even pages carry mirrored
% margins and the text block visibly steps left and right on screen. Passing
% oneside to \documentclass does not help, because the class applies its own
% twoside afterwards and the later declaration wins. Average the two instead.
\AtBeginDocument{%
  \addtolength{\oddsidemargin}{\evensidemargin}%
  \setlength{\oddsidemargin}{0.5\oddsidemargin}%
  \setlength{\evensidemargin}{\oddsidemargin}%
}

% \bmhead is a level five heading while the class sets bookmarksdepth=5, so
% the back matter jumps from level one straight to level five and hyperref
% reports a bookmark level warning once per heading.
\hypersetup{bookmarksdepth=3}

% Reference entries contain unbreakable strings, so TeX sets the surrounding
% lines loose. \sloppy alone permits the loose line but leaves \hbadness at
% 1000, so both are required, applied at \begin{document} to land after
% natbib's own hooks.
\AtBeginDocument{%
  \let\snOrigThebibliography\thebibliography
  \renewcommand{\thebibliography}[1]{%
    \snOrigThebibliography{#1}\sloppy\hbadness=10000\relax}%
}

\begin{document}

\title[Wavelet Denoising and Transferred Architectures for Cough Detection]{Wavelet
Denoising Before Spectrogram Extraction Improves Cough Detection for Moderate
Capacity Networks but Not for Densely Connected Ones}

\author[1]{\fnm{Akhilesh Bharadhwaj} \sur{S}}\email{akhilesh.bharadhwaj2024@vitstudent.ac.in}

\author[1]{\fnm{Vishnuvarthan} \sur{M}}\email{vishnuvarthan.m2024@vitstudent.ac.in}

\author[1]{\fnm{Nithin} \sur{Kumar}}\email{nithinkumar.j2024a@vitstudent.ac.in}

\author[1]{\fnm{Ajay} \sur{Pokkali}}\email{ajay.pokkali2024@vitstudent.ac.in}

\author*[1]{\fnm{Ravi} \sur{Tiwari}}\email{ravi.tiwari@vit.ac.in}

\affil*[1]{\orgdiv{School of Electronics Engineering (SENSE)},
  \orgname{Vellore Institute of Technology},
  \orgaddress{\city{Chennai}, \state{Tamil Nadu}, \country{India}}}

\abstract{Automatic detection of coughs in continuously recorded audio is the
gating step for any acoustic respiratory screening system, and its reliability
determines how often such a system raises an alarm on material that contains no
cough at all. We study that step on crowdsourced recordings, where background
noise varies without constraint, and we examine three design choices
together: whether the network receives a linear or a perceptually scaled
spectrogram, whether the waveform is denoised with a wavelet transform before
the spectrogram is formed, and which convolutional architecture performs the
classification. Four models are evaluated, a compact network trained from
scratch and three architectures transferred from a large image corpus, each
with and without denoising, giving eight configurations measured on an
identical partition of 4,577 held out recordings. Perceptual frequency scaling
improves every architecture it is applied to. Denoising helps the two models
of moderate capacity, raising the best accuracy to 96.53 percent with a
precision of 97.59 percent and a recall of 97.31 percent, but it does not help
the densely connected network, whose accuracy is unchanged within rounding. We
also recover the full error structure of every configuration from the reported
rates and show that the class imbalance of the corpus conceals a wide spread
in specificity, from 90.84 to 95.73 percent, that plain accuracy does not
reveal. Ranking the configurations by balanced accuracy preserves the
conclusion while narrowing the apparent margins considerably, and we report
the reordering this causes because it changes which model a deployment should
prefer.}

\keywords{Respiratory screening, Acoustic event detection, Crowdsourced audio,
Perceptual frequency scaling, Class imbalance, Signal preprocessing}

\maketitle

\section{Introduction}\label{sec:intro}

A cough is one of the few symptoms that produces a signal a microphone can
capture without contact, without a clinician present and without the person
being examined doing anything beyond carrying a telephone. That property has
made acoustic analysis of coughing an attractive route to population scale
respiratory screening, and the interest intensified during the COVID-19
pandemic, when the need to triage large numbers of people quickly exceeded the
capacity of testing programmes in most countries. Systems built in that period
sought to infer disease status directly from a recorded cough, and several
reported encouraging results on data gathered through web forms and mobile
applications~\cite{brown2020, imran2020, laguarta2020}. The conditions of the
pandemic also produced the open corpora that made comparison across research
groups possible, of which the COUGHVID collection is among the
largest~\cite{orlandic2021}.

Beneath the question of what a cough reveals lies a prior question that is
easier to state and harder to dismiss: is the sound a cough at all. Every
screening pipeline that operates on continuous audio must first decide which
segments contain coughing and which contain speech, laughter, door closures,
traffic, music or silence. That decision is the gate through which all later
inference passes. A missed cough removes a person from consideration entirely,
and a false detection sends noncough audio into a disease classifier that was
never trained on it, producing an output with no meaning that is nonetheless
reported with a confidence value. In a screening context the second error is
the more insidious, because it generates referrals that consume clinical time
and erode confidence in the system among the people who have to act on its
output. Detection therefore deserves to be studied as a problem in its own
right rather than as a preprocessing detail, and that is the subject of this
paper.

The acoustic properties of a cough make the detection problem tractable in
principle. A cough is a forced expiratory manoeuvre with a characteristic
three phase structure: an explosive onset as the glottis opens against raised
subglottal pressure, an intermediate phase of rapid airflow, and a voiced or
unvoiced termination. The onset produces a broadband transient with energy
extending well above the range occupied by conversational speech, and the
sequence of phases gives a temporal envelope that is distinctive even when the
spectrum alone is ambiguous. Early systems exploited exactly these properties
through features designed by hand, including Mel frequency cepstral
coefficients, spectral centroid, zero crossing rate and short term energy,
often with a classifier of modest capacity on
top~\cite{swarnkar2013, hoyosbarcelo2018, mongealvarez2019}. Such features are
interpretable and inexpensive, and they remain appropriate where computation is
severely constrained, but they encode a fixed prior about what distinguishes a
cough and they generalise poorly when recording conditions depart from those
the designer anticipated.

Learned representations displaced designed ones for the same reasons they did
elsewhere in audio processing. Presenting a time frequency image to a
convolutional network allows the network to discover which regions of the
plane carry discriminative information, and work on general audio event
recognition established both that this approach is effective and that
architectures developed for natural images transfer to it with little
modification~\cite{piczak2015, hershey2017, salamon2017}. The transfer is
possible because a spectrogram, once rendered as an image, exhibits the local
structure that convolutional kernels are suited to, and because large image
corpora supply initial filters that are generic enough to be useful on inputs
quite unlike photographs~\cite{deng2009, pan2010}. The practical consequence is
that a cough detector can be built from an architecture trained on millions of
images, retrained on a corpus of a size that a research group can actually
assemble.

Two choices within that recipe are usually made without being examined.
The first concerns the frequency axis of the spectrogram. A short time
Fourier transform spaces its bins linearly, which allocates the same
resolution to the region between 7 and 8~kHz as to the region between 100 and
200~Hz, although the latter carries far more of what distinguishes one
respiratory sound from another. Mapping the axis onto a perceptual scale
redistributes that resolution, compressing the upper range and expanding the
lower, following the relationship established in early work on speech
features~\cite{davis1980}. The second concerns noise. Crowdsourced recordings
are made on unknown equipment in uncontrolled surroundings, so the corpus
contains traffic, electrical hum, wind, television audio and clipping, none of
it labelled. Wavelet decomposition offers a principled way to attenuate such
interference before any spectrogram is formed, because it represents a signal
at multiple resolutions simultaneously and allows the components that carry
broadband transient energy to be separated from those that do
not~\cite{mallat1989}, with thresholding rules that have a firm statistical
basis~\cite{donoho1994, donoho1995}.

Whether denoising is worth performing is less obvious than it first appears,
and the reason is that it interacts with the capacity of the network placed
after it. A deep network with extensive feature reuse may already learn
representations that ignore interference, in which case removing that
interference beforehand buys nothing and risks discarding genuine high
frequency content that the network had learned to exploit. A network of more
modest capacity, or one trained for fewer epochs, may lack that ability and
benefit substantially. This interaction is rarely tested, because most studies
fix one architecture and vary the preprocessing, or fix the preprocessing and
vary the architecture, so the question of whether the benefit of denoising
depends on what follows it remains open. The present work tests precisely that
interaction by crossing both factors on one corpus with one partition.

This paper makes the following contributions.

\begin{enumerate}
\item We evaluate eight configurations on an identical partition of the
COUGHVID corpus, crossing two spectrogram representations and the presence or
absence of wavelet denoising across four convolutional architectures, so that
the effect of each factor can be read without the confound of a differing
split or a differing training budget.
\item We show that the benefit of wavelet denoising depends on the
architecture that follows it. The two models of moderate capacity gain, by
0.77 and 0.35 accuracy points respectively, whereas the densely connected
network does not, its accuracy changing by less than one tenth of a point in
the opposite direction.
\item We recover the complete confusion structure of all eight configurations
from the published rates and the composition of the test partition, and report
specificity and balanced accuracy, neither of which the accuracy figure alone
exposes on a corpus whose classes stand in a ratio greater than two to one.
\item We show that ranking by balanced accuracy preserves the identity of the
best configuration but compresses the margins between the remainder and
reorders the models in the middle of the table, and we set out what this
implies for a deployment that must choose between a stronger detector and a
smaller one.
\end{enumerate}

The remainder of the paper is organised as follows.
Section~\ref{sec:related} reviews the relevant literature and identifies the
gap addressed here. Section~\ref{sec:design} specifies the processing pipeline
and the four architectures. Section~\ref{sec:setup} describes the corpus, the
training configuration, the metrics and the information needed to repeat the
work. Section~\ref{sec:results} presents the results, the recovered error
structure and the limitations of the study. Section~\ref{sec:conclusion}
concludes.

\section{Related work and its limitations}\label{sec:related}

This section reviews five strands of work and states what each leaves
unresolved. Section~\ref{sec:rel-handcrafted} covers detection from features
designed by hand. Section~\ref{sec:rel-corpora} covers the open corpora that
made comparative evaluation possible and the screening systems built on them.
Section~\ref{sec:rel-transfer} covers spectrogram inputs and transferred
architectures. Section~\ref{sec:rel-eval} covers evaluation practice in this
literature and why reported figures are difficult to compare across studies.
Section~\ref{sec:rel-denoise} covers denoising for respiratory audio and
states the gap this study addresses.

\subsection{Detection from designed features}\label{sec:rel-handcrafted}

The first generation of automatic cough detectors described the signal through
a small set of quantities chosen for their physical interpretation and passed
them to a classifier of modest capacity. Swarnkar and colleagues separated wet
from dry coughs in paediatric patients using features of this kind, showing
that the distinction clinicians draw by ear has a measurable acoustic
correlate~\cite{swarnkar2013}. Hoyos-Barcel\'o and colleagues addressed the
engineering problem directly, implementing a nearest neighbour classifier
efficiently enough to detect cough events on a smartphone in real
time~\cite{hoyosbarcelo2018}, and the same group later described a detector
built on local moments computed over the time frequency
plane~\cite{mongealvarez2019}. Nemati and colleagues considered the broader
task of distinguishing cough types rather than merely detecting
them~\cite{nemati2019}.

The features themselves were chosen to reflect properties of the signal that
can be described in words. Cepstral coefficients summarise the gross shape of
the spectral envelope in a small number of values and so capture the resonant
structure imposed by the airway. Spectral centroid locates the centre of mass
of the spectrum and separates sounds that are bright from those that are
dull. Zero crossing rate responds to the proportion of high frequency content
and distinguishes the noisy portion of a cough from voiced material. Short
term energy traces the envelope and therefore the phase structure that gives
a cough its characteristic shape in time. A classifier presented with these
quantities is working with a compact description that a clinician can be
shown and, in principle, dispute.

These systems share two strengths that are easy to undervalue. They are cheap
enough to run continuously on battery powered hardware, and their decisions
can be traced to quantities a clinician can be shown. Their weakness is that
the feature set encodes an assumption about what a cough sounds like, fixed
before any data is seen. Where the recording conditions match that assumption
the systems perform well; where they do not, performance falls in ways that
cannot be corrected by retraining the classifier alone, because the
information the classifier would need has already been discarded by the
feature extractor. Crowdsourced audio departs from any fixed assumption by
construction, which is what motivated the move to learned representations.

\subsection{Open corpora and screening systems}\label{sec:rel-corpora}

Progress in this area was limited for a long period by the absence of data
that different groups could share. The pandemic changed that. Brown and
colleagues assembled crowdsourced respiratory recordings and showed that
simple classifiers applied to them could separate reported COVID-19 cases from
controls at rates well above chance~\cite{brown2020}. Sharma and colleagues
released the Coswara database, covering breathing, cough and voice
together~\cite{sharma2020}. Orlandic and colleagues released COUGHVID, the
corpus used here, and documented both its scale and the variability that makes
it demanding: the recordings arrive through a web interface from unknown
devices in unknown surroundings, and a substantial fraction contain no cough at
all~\cite{orlandic2021}.

Screening systems followed quickly. Imran and colleagues built a deployable
application around a deep classifier operating on cough
audio~\cite{imran2020}, Laguarta and colleagues reported that a model trained
on cough recordings alone could discriminate COVID-19 status~\cite{laguarta2020},
Pahar and colleagues examined the same question across smartphone recordings
gathered globally~\cite{pahar2021}, and Coppock and colleagues trained a
network end to end on breath and cough audio~\cite{coppock2021}. Later work
pursued architectural refinement for the same task~\cite{dentamaro2022,
vijayakumar2021, lella2021}, and independent groups revisited the crowdsourced
data to test how far the early claims
generalised~\cite{son2022}. Ijaz and colleagues survey the field as a
whole~\cite{ijaz2022}.

A reckoning followed the first wave of results. Several of the early claims
rested on corpora in which the positive and negative groups differed in more
than their disease status, because the people who contributed them were
recruited through different channels, recorded in different circumstances and
sometimes on different classes of device. A model can separate such groups
without learning anything about respiratory disease, and independent
re examination of crowdsourced material has reported figures below those
initially announced once such confounds are controlled~\cite{son2022}. The
episode left the field with a sharper sense of what an acoustic screening
claim requires: a negative class drawn from the same population and the same
recording conditions as the positive one, a partition that separates
contributors rather than recordings, and a baseline that shows what the
prior alone would achieve.

That history bears directly on the detection problem studied here, because a
detector is the component that decides which audio a disease model ever sees.
If detection errors correlate with recording conditions, and crowdsourced
recording conditions correlate with contributor, then a detector can
introduce exactly the kind of confound the field has learned to guard
against, before the disease model is reached at all.

The limitation relevant here is one of emphasis rather than of method. Almost
all of this work takes the presence of a cough as given and asks what the
cough indicates. The recordings are segmented, often manually or with a
detector whose own error rate is not reported, and the screening model is
evaluated on the segments that survive. Any error introduced at that stage is
therefore absent from the reported figures while remaining present in any
deployment, and since the detection stage is where noncough audio is rejected,
its false positive rate propagates directly into the referral rate of the
complete system.

\subsection{Spectrogram inputs and transferred architectures}\label{sec:rel-transfer}

Treating a time frequency representation as an image allows a convolutional
network to learn its own features, and the approach was established on general
audio before it was applied to respiratory sounds. Piczak showed that a
convolutional network on spectrogram inputs is competitive for environmental
sound classification~\cite{piczak2015}, Salamon and Bello improved the result
with augmentation designed for audio~\cite{salamon2017}, and Hershey and
colleagues examined how image architectures behave at large scale on
audio~\cite{hershey2017}, supported by the release of a correspondingly large
labelled collection~\cite{gemmeke2017}.

The architectures themselves originate in image recognition. Residual
connections allowed networks to be made substantially deeper without the
degradation that had previously accompanied depth~\cite{he2016}. Dense
connectivity, in which every layer receives the feature maps of all preceding
layers, strengthens gradient flow further and encourages features to be reused
rather than recomputed~\cite{huang2017}. Separately, work on efficiency
produced architectures intended for devices with limited memory and power, of
which the inverted residual design with linear bottlenecks is a widely adopted
example~\cite{sandler2018}. All three transfer to spectrogram inputs by
replacing the classification head, with initial weights taken from a large
image corpus~\cite{deng2009}; the general conditions under which such transfer
succeeds are surveyed by Pan and Yang~\cite{pan2010}.

Why transfer from photographs to spectrograms works at all deserves a word,
because the two kinds of image are not alike. A photograph is approximately
stationary under translation in either direction, so a feature is equally
meaningful wherever it appears. A spectrogram is not: translating a pattern
along the time axis leaves its meaning intact, whereas translating it along
the frequency axis changes what it represents. A convolution applied over both
axes therefore imposes an invariance on the frequency axis that the signal
does not possess. Transfer nonetheless succeeds, and the reason appears to be
that the early layers of an image network learn generic detectors of edges,
corners and oriented texture, and these remain useful on a spectrogram because
the structures that matter there, harmonic ridges, broadband vertical
transients and the boundaries of energy bursts, are themselves edges and
textures. The layers whose features are specific to natural images lie deeper,
and those are the layers that retraining adjusts.

What is less often examined is the frequency axis these architectures receive.
A linear axis and a perceptually scaled one present the same recording as
visibly different images, and the convolution kernel, which is spatially
uniform, therefore sees a different problem in each case. The perceptual
mapping concentrates resolution where respiratory sounds differ most, and the
relationship used for it was established in the speech recognition
literature~\cite{davis1980}. The choice also interacts with the resize that
precedes a transferred network, since an architecture expecting a square
input of fixed size will compress or stretch whatever representation it is
given, and the amount of distortion this introduces depends on how many
frequency bins the representation carried in the first place. Studies
comparing the two axes on cough detection specifically, with everything else
held constant, are uncommon.

\subsection{Evaluation practice in this literature}\label{sec:rel-eval}

Reported figures in cough analysis are difficult to compare across studies,
and the reasons are worth setting out because they shape how the present
results should be read.

The first reason is that the corpora differ in what they contain. A corpus
assembled from clinical recordings contains coughs produced by people who are
unwell, captured on consistent equipment at a controlled distance. A
crowdsourced corpus contains coughs produced on request by a largely healthy
population on whatever device they happened to hold, at distances and levels
that vary by orders of magnitude, with background content that ranges from
silence to television audio. Detection is a different problem in the two
settings, and a figure obtained in one does not transfer to the
other~\cite{orlandic2021, ijaz2022}. The same applies to the negative class.
What counts as a noncough recording is defined differently by different
collections, and a detector evaluated against a negative class consisting
mostly of speech will report a different error rate from one evaluated
against a negative class dominated by mechanical noise.

The second reason is the partition. Many reported results use a random split
at the level of the recording, which allows multiple recordings from one
contributor to appear on both sides of the division. Where contributors
differ systematically in device, room and vocal character, such a split lets
a model recognise the contributor rather than the event, and the resulting
figure overstates what the model would achieve on an unseen person. Splitting
by contributor avoids this but is possible only where contributor identifiers
are retained, which not every collection provides.

The third reason is the choice of summary measure. Cough corpora are commonly
imbalanced, and on an imbalanced partition accuracy is dominated by the
majority class, so two systems with identical accuracy can differ
substantially in how they distribute their errors. Precision, recall and
$F_1$ expose that distribution, which is why all four are reported here,
and the balance between them is the quantity a deployment actually
selects on. Comparable studies vary in which subset they report, and a study
reporting accuracy alone leaves a reader unable to tell whether an
improvement came from detecting more coughs or from rejecting more noncough
material~\cite{son2022}.

The consequence for the present work is methodological rather than
substantive. Because cross study comparison is unreliable for these reasons,
the design adopted here holds the corpus, the partition, the optimiser and the
evaluation fixed and varies only the two factors under examination. The
comparisons that carry the argument are therefore internal to this study, and
we draw no conclusion from the numerical proximity of our figures to those
reported elsewhere.

\subsection{Denoising for respiratory audio, and the gap}\label{sec:rel-denoise}

Wavelet analysis is well matched to signals whose informative content is
transient. Decomposing a waveform across scales isolates abrupt broadband
events, such as the explosive onset of a cough, from interference that is
either narrowband or slowly varying, and the multiresolution framework that
makes this possible is due to Mallat~\cite{mallat1989}. The statistical basis
for discarding coefficients was established by Donoho and Johnstone, who showed
that thresholding in a wavelet basis is near optimal for recovering a signal
from additive noise~\cite{donoho1994} and characterised the behaviour of soft
thresholding in particular~\cite{donoho1995}. Applied to a cough recording, the
approach attenuates the background while leaving the transient structure that
carries the discriminative information largely intact.

Reported results on respiratory audio generally favour denoising, but the
evidence has a consistent shape that limits what can be drawn from it. A study
fixes one architecture, applies denoising, observes an improvement and
concludes that denoising is beneficial. This establishes that denoising helps
that architecture on that corpus. It does not establish that the benefit is a
property of denoising rather than of the pairing, and there is a specific
reason to doubt that it is. Denoising and architectural capacity address the
same problem by different means. Removing interference before the network sees
it, and building a network able to ignore interference, are substitutes rather
than complements, so a network that already achieves invariance through its own
structure should gain little from preprocessing and may lose whatever genuine
high frequency content the preprocessing removes along with the noise.

The gap this study addresses follows directly. We hold the corpus, the
partition, the optimiser and the evaluation fixed, and vary the representation
and the denoising across four architectures that differ substantially in depth
and in the mechanism by which they reuse features. This design allows the
question to be posed as an interaction rather than as a main effect, and
Section~\ref{sec:results} shows that the interaction is real: the benefit of
denoising is not a constant of the method but depends on what is placed after
it.
